\section{Conclusion}
\label{sec:conclusion}

In this paper, we analyzed an existed reversible data hiding method proposed by Yin et al., and we proposed three improvements, within which, two of them have significant incremental in embedding capacity that has been shown in Section~\ref{sec:experimental}, while the last improvement aimed to make the method more stable and generalize that can accept cover images of larger size. 

Since the method within our discussion is a reversible data hiding method, 
which can losslessly restore the cover image after extracting secret message, 
and also because the stego image is indeed an encrypted image, 
we omitted the comparison using PSNR and MSE in Section~\ref{sec:experimental} for their value will always be $+\infty$ and $0$.
We focused our comparison on the embedding capacity, i.e., BPP, and the experimental result of applying both method shows that our proposal is superior even in the worst case scenario when applying on rough images. 